\begingroup
\fontsize{9.5pt}{11pt}\selectfont
\setlength{\tabcolsep}{3pt}
\begin{longtable}{@{}p{0.11\linewidth}p{0.12\linewidth}p{0.12\linewidth}p{0.07\linewidth}p{0.26\linewidth}p{0.26\linewidth}@{}}
\caption{Target-by-axis qualification audit corresponding to main-text Figure 7. Each row links a target-specific evidence state to the audited claim source, states the current evidence, and specifies what evidence would be required to change that state. Not evaluated and not applicable are not negative results.}
\label{tab:supp-evidence-axis-audit}\\
\toprule
Target & Evidence axis & State & Claim & Current evidence & Next evidence needed \\
\midrule
\endfirsthead
\toprule
Target & Evidence axis & State & Claim & Current evidence & Next evidence needed \\
\midrule
\endhead
Grade and phenotype & Frozen primary & supported & C02 & Source-cohort grade and phenotype signals were supported. & Retain the locked source probes and analysis units. \\
Grade and phenotype & Cross-cohort transport & supported & C02 & Direction was retained without refitting in PANDA and was concordant with PRECISE grade evaluation. & Add patient-level uncertainty where source data permit and test further institutions. \\
Grade and phenotype & Clinical increment & not evaluated & C02 & No incremental clinical-covariate claim was defined for these morphology-linked targets. & Define a clinically meaningful reference model before making an increment claim. \\
Grade and phenotype & Site behavior & supported & C02 & Transfer directions were retained at both PANDA institutions; this was not a causal site analysis. & Use site-specific patient-level intervals and acquisition metadata in additional cohorts. \\
Grade and phenotype & Setting stability & supported & C02;C08 & Correlated setting audits retained the transported directions. & Confirm stability on independently sampled patients rather than reused settings. \\
Grade and phenotype & Endpoint fidelity & not applicable & C02 & No time-to-event endpoint was used for these targets. & Not applicable to the present grade and phenotype claims. \\
PTEN & Frozen primary & supported & C03 & The within-TCGA frozen-primary association was above chance. & Repeat the locked evaluation in an external cohort with compatible PTEN labels. \\
PTEN & Cross-cohort transport & not evaluated & C03 & No compatible external PTEN cohort was available in this study. & Acquire an external cohort with target-matched PTEN assays. \\
PTEN & Clinical increment & unresolved & C03 & Held-out increment beyond grade was not established. & Test a locked image increment against grade and fuller clinical covariates in new patients. \\
PTEN & Site behavior & not evaluated & C03 & The manuscript does not establish external site transport for PTEN. & Use a multisite external design with site and acquisition metadata. \\
PTEN & Setting stability & supported & C03;C08 & The within-cohort direction remained above chance across the evaluated grid. & Confirm the direction under independent cohort resampling. \\
PTEN & Endpoint fidelity & not applicable & C03 & PTEN loss was a binary molecular target rather than a clinical time-to-event endpoint. & Not applicable to the present PTEN claim. \\
SPOP & Frozen primary & unsupported in frozen design & C04 & The frozen primary configuration did not support the proposed effect. & Increase target-appropriate sampling without changing the frozen primary interpretation. \\
SPOP & Cross-cohort transport & not evaluated & C04 & No compatible external SPOP cohort was evaluated. & Test a locked probe in an external cohort with mutation labels and tumor-content control. \\
SPOP & Clinical increment & not evaluated & C04 & A grade-independent SPOP increment was not established in the approved analysis set. & Evaluate incremental value with patient-disjoint nested analysis if scientifically justified. \\
SPOP & Site behavior & unresolved & C04 & Defined site intervals included or touched chance and one site was single-class. & Collect larger site-balanced data with assay and acquisition metadata. \\
SPOP & Setting stability & context-sensitive & C04;C08 & The result crossed the null across multiple correlated settings. & Resolve tumor sampling and representation sensitivity in independent data. \\
SPOP & Endpoint fidelity & not applicable & C04 & SPOP mutation was a binary molecular target rather than a clinical time-to-event endpoint. & Not applicable to the present SPOP claim. \\
Androgen-receptor activity & Frozen primary & supported & C05 & The pooled frozen-primary direction was positive. & Retain the pooled result as within-cohort evidence only. \\
Androgen-receptor activity & Cross-cohort transport & not evaluated & C05 & No compatible external androgen-receptor activity cohort was evaluated. & Use an external cohort with a compatible activity assay. \\
Androgen-receptor activity & Clinical increment & unresolved & C05 & Held-out grade-independent increment was not established. & Test a locked incremental model against grade and fuller clinical structure. \\
Androgen-receptor activity & Site behavior & unresolved & C05 & Every stored site-specific interval crossed zero and acquisition effects could not be separated. & Disentangle site from scanner stain preparation and case mix. \\
Androgen-receptor activity & Setting stability & supported & C05;C08 & The positive direction was retained across the correlated setting grid. & Confirm with independently sampled sites and patients. \\
Androgen-receptor activity & Endpoint fidelity & not applicable & C05 & Androgen-receptor activity was a continuous molecular target rather than a time-to-event endpoint. & Not applicable to the present activity-score claim. \\
Recurrence & Frozen primary & context-sensitive & C06 & The transferred post-hoc risk signal depended on representation and endpoint. & Lock the source model before a new independent target-cohort evaluation. \\
Recurrence & Cross-cohort transport & context-sensitive & C06 & LEOPARD-to-TCGA transfer was present in some settings but was not endpoint-invariant. & Test the locked score in an external cohort with a harmonized endpoint. \\
Recurrence & Clinical increment & unresolved & C07 & Full clinical and site comparisons did not establish robust incremental value. & Use an adequately powered same-patient comparison against a clinical reference locked before evaluation. \\
Recurrence & Site behavior & not evaluated & C06;C07 & Site was included in fuller models but site-specific recurrence transport was not isolated. & Evaluate site transport with sufficient events per site. \\
Recurrence & Setting stability & context-sensitive & C06;C08 & Effect size and direction depended on encoder scale and related settings. & Use nested locked model selection and independent validation. \\
Recurrence & Endpoint fidelity & context-sensitive & C06;C07 & Reconstructed recurrence and official PFI produced different evidentiary conclusions. & Harmonize endpoint definitions and retain endpoint-specific reporting. \\
\bottomrule
\end{longtable}
\endgroup
